\section{Knowledge of Capital Markets}

\medskip

\subsection{Market Participants}

\medskip

One of the primary purposes of the securities industry is to match investors, who have money, with issuers that need money.

\medskip

\begin{definition}[Issuer]
An issuer is a legal entity that sells securities in order to finance its operations. Examples of issuers are government agencies, governments themselves, corporations, and banks. Issuers primarily issue either debt securities (bonds) or equity securities (stocks). 
\end{definition}

Debt securities are publicly traded loans, referred to as bonds, notes or debt instruments. A bond represents the amount of indebtedness (principal) that the issuer owes to the investor. The investor (purchaser) is considered a \textit{creditor} and has lent their money to the issuer for the specified period (till maturity). The bond issuer is expected to pay the principal balance at the end of maturity, as well as possible interest payments. If any payment is missed then the bond is in \textit{default}.

Equity securities represent ownership interest of a corporation, and are issued by corporations to raise capital. If the company is profitable then those with ownership interest may be entitled to profits (dividends). This security does not typically have a maturity date and dividend payment is voluntary for the issuer. 

Financial firms connect issuers and investors, and are broadly categorized into either \textit{broker-dealers} or \textit{investment advisers}.

\bigskip

\subsubsection{Broker-Dealers}

\medskip

A broker engages in the business of \textit{agency} transactions in securities for the account of others. In an agency transaction, the broker acts strictly as an intermediary connecting buyers and sellers without taking ownership of the securities into inventory, thereby avoiding principal risk and earning compensation through commissions rather than markups.

A dealer is a person that engages in the business of buying or selling securities for its own account. Dealers engage in \textit{principal} transactions, where securities are bought directly from clients and held in inventory, or sold to clients from inventory. 

Since firms can continuously execute trades as either a broker or a dealer, it is convenient to refer to the whole lot as broker-dealers. 

\bigskip

\subsubsection{Broker-Dealer Departments}

\medskip

Brokerage firms (broker-dealers) are divided into many distint departments. The most common of which are:
\vspace{-1.5ex}
\begin{itemize}[label=$\circ$, itemsep=0.6em]
    \item Investment Banking (IB): IB is an area that works directly with issuers to arrange and structure their securities offerings. Investment bankers are often referred to as \textit{underwriters} of securities, and may also assist companies with mergers or acquisitions (M\&A) or restructuring due to bankruptcy. 
    \item Research: Analysts study the market and individual issuers to make recommendations (typically buy, sell, or hold).
    \item Sales: Market individual stocks or bonds (or packaged products such as mutual funds) to retail investors and institutions. Historically referred to as stock or bond brokers, but more recently as \textit{registered representatives (RRs)} or \textit{investment advisor representatives (IARs)}.
    \item Trading: Handles the execution of trades for the firm's clients and the firm's own proprietary account. Trades can occur in electronic marketplaces (e.g. Nasdaq) or hybrid marketplaces (e.g., NYSE). 
    \item Operations: Ensures all paperwork, funds, and securities transfers associated with trades or processing is handled efficiently and in compliance with industry standards and regulations. Functions may include generating customer statements, confirmations, tax records, and assisting in transfer of funds or securities. Also responsible for ensuring that client and firm assets are organized and safeguarded. 
\end{itemize}

\bigskip

\subsubsection{Market Makers}

\medskip

Market makers post quotes (two-sided) for buyings and/or selling securities at specific prices. Often, market makers are \textit{required} to make regular bids (buy requests) and offers (sell offerings) and to meet specific capital requirements. For example, where the firm is willing to buy 1,000 shares at 17.05 and/or sell 2,000 shares at 17.08:

\begin{table}[htbp]
  \centering
  \begin{tabularx}{\linewidth}{CCC}
    \midrule
    \textbf{Bid} & \textbf{Ask} & \textbf{Size} \\[0.5em]
    \rowcolor{valuerow}
    \rule[-9pt]{0pt}{24pt}17.05 & 17.08 & $1{,}000 \times 2{,}000$ \\
    \addlinespace[6pt]
    \makecell{Price at which\\ the firm will buy} &
    \makecell{Price at which\\ the firm will sell} & \\
    \midrule
  \end{tabularx}
\end{table}

